\documentclass[11pt,twocolumn,a4paper]{article}

% ─── Packages ─────────────────────────────────────────────────────────
\usepackage[margin=2.0cm,top=2.4cm,bottom=2.4cm]{geometry}
\usepackage[utf8]{inputenc}
\usepackage[T1]{fontenc}
\usepackage{lmodern}
\usepackage{microtype}
\usepackage{amsmath,amssymb,amsthm,mathtools}
\usepackage{bm}
\usepackage{graphicx}
\usepackage{booktabs}
\usepackage{array}
\usepackage{enumitem}
\usepackage{float}
\usepackage{caption}
\usepackage{natbib}
\usepackage{xcolor}
\usepackage{listings}
\usepackage[colorlinks=true,linkcolor=blue!60!black,
            citecolor=blue!60!black,urlcolor=blue!60!black]{hyperref}
\usepackage{fancyhdr}
\usepackage{titlesec}
\usepackage{cleveref}

% ─── Page Style ───────────────────────────────────────────────────────
\pagestyle{fancy}
\fancyhf{}
\fancyhead[L]{\small\textit{Mishima: Recalling Authored Work by Propagation}}
\fancyhead[R]{\small\thepage}
\renewcommand{\headrulewidth}{0.4pt}

\titleformat{\section}{\large\bfseries}{\thesection.}{0.5em}{}
\titleformat{\subsection}{\normalsize\bfseries}{\thesubsection.}{0.5em}{}

% ─── Theorem Environments ─────────────────────────────────────────────
\newtheorem{theorem}{Theorem}[section]
\newtheorem{lemma}[theorem]{Lemma}
\newtheorem{corollary}[theorem]{Corollary}
\newtheorem{proposition}[theorem]{Proposition}
\theoremstyle{definition}
\newtheorem{definition}[theorem]{Definition}
\newtheorem{example}[theorem]{Example}
\newtheorem{principle}[theorem]{Principle}
\theoremstyle{remark}
\newtheorem{remark}[theorem]{Remark}

% ─── Custom Commands ──────────────────────────────────────────────────
\newcommand{\R}{\mathbb{R}}
\newcommand{\N}{\mathbb{N}}
\newcommand{\Rpos}{\R_{>0}}
\newcommand{\flo}{\beta}
\newcommand{\med}{\mathcal{M}}
\newcommand{\Gph}{\Gamma}
\newcommand{\cut}{\operatorname{cut}}
\newcommand{\wt}{w}
\newcommand{\res}{\varrho}
\newcommand{\pot}{\kappa}
\newcommand{\aln}{\sigma}
\newcommand{\rec}{\mu}
\newcommand{\Om}{\Omega}
\newcommand{\Reached}{\mathcal{C}}

% ─── Listings ─────────────────────────────────────────────────────────
\lstdefinestyle{mma}{
  basicstyle=\ttfamily\footnotesize,
  keywordstyle=\bfseries,
  commentstyle=\itshape\color{gray!80!black},
  showstringspaces=false,
  breaklines=true,
  columns=fullflexible,
  frame=single,
  framesep=4pt,
  rulecolor=\color{gray!50},
  morekeywords={floor,module,probe,commit,seek,not,toward,via,until,
                closure,converge,otherwise,decline,yield,ladder,rung,at,
                observe,assert,emit,let,as,when}
}
\lstset{style=mma}

\captionsetup{font=small,labelfont=bf}

% ══════════════════════════════════════════════════════════════════════
\begin{document}

\twocolumn[
\begin{@twocolumnfalse}
\begin{center}

{\LARGE\bfseries Mishima: Recalling Authored Work by Propagation\\[4pt]
A Language in which a Search States What It Excludes,\\
and Ends in Closure or Honest Decline}

\vspace{12pt}

{\large Kundai Farai Sachikonye}

\vspace{6pt}

{\small\itshape Independent Researcher \\[0.3em]
\texttt{kundai.sachikonye@wzw.tum.de}}

\vspace{6pt}{\small August 2026}

\vspace{16pt}

\begin{minipage}{0.92\textwidth}
\textbf{Abstract.}
We specify \emph{mishima}, a language for asking questions of an accumulated
record of music production. Its subject is the request a producer cannot
currently make of any tool: \emph{find the thing I made before, the one that
sounded like this and not like those, and tell me what was essential to it}. We
argue that such a request is not vague but negatively specified, and we make the
negation syntactically obligatory: a search that does not state what it excludes
is rejected at parse time, because a region individuated against nothing is not a
region. The language's cost model counts commitments rather than instructions, so
that exploring an enormous space while committing rarely is inexpensive however
long the exploration. We prove that a single distinction cannot be sharpened
without bound---the cost of thinning a separator diverges while its return is
bounded---and therefore that fine resolution must be reached by composing
ordinary distinctions in series, a \emph{ladder} whose composite power is
$1-\prod_i(1-\pot_i)$ and which saturates exactly when $\sum_i \pot_i$ diverges.
Measured over four thousand rungs, a divergent power sequence drives the residual
gap to $2.5\times10^{-4}$ while a convergent one leaves it at $0.29$, a
separation of three orders of magnitude. We give a termination criterion---%
\emph{closure}, that every uninvoked probe lands where the search has already
been---and prove it strictly stronger than any confidence threshold; over five
registries a fixed threshold was satisfied after a single probe in every case
while closure required the whole registry, and every registry terminated
contested. Contested termination is a typed value carrying the irreconcilable
readings found, not an error, because a system holding no expectation cannot
adjudicate between them. We give the complete lexical and context-free grammar
for the \texttt{.mma} file, a type system carrying resolution and power in the
type, a small-step semantics, and static diagnostics that reject an incoherent
ladder before it runs. A prototype front end and sixteen deterministic
experiments accompany the paper. Nothing outside this paper is required to read
it.
\end{minipage}

\vspace{8pt}
\noindent\textbf{Keywords:}
domain-specific language, music production, negation, resolution ladder,
closure, honest decline, water-filling, type system, provenance

\vspace{18pt}
\end{center}
\end{@twocolumnfalse}
]

%──────────────────────────────────────────────────────────────────────
\section{Introduction}
\label{sec:intro}
%──────────────────────────────────────────────────────────────────────

\subsection{The question a producer cannot ask}

Suppose a producer wants the bass sound they made two years ago. Not a similar
one---that one, or something that extends it. What can they do?

They can open old sessions and listen. This works and is what everyone does, and
it scales as badly as it sounds. They can search filenames, which requires that
the name recorded the thing they now want, which it did not, because at the time
they did not know they would want it. They can send a bounce to a similarity
service and receive back other people's records that resemble it, which is a
different question.

None of these is a failure of tooling to be fixed by a better index. The
difficulty is prior: \emph{the material does not record what was decided}. A
bounce is the output of a chain of decisions with every trace of the decisions
removed; a session file is the terminal state of those decisions with the
sequence removed. Neither contains the relation between what was tried and what
was kept, and it is the relation that a later question is about.

Mishima assumes that relation has been recorded---by a runtime specified
separately, and summarised self-containedly in \S\ref{sec:substrate}---and is the
language for asking questions of it. This paper is about the asking.

\subsection{Why the question is negatively specified}

``The sound I am looking for'' is routinely described as vague. We think this is
a mistake about the grammar of such requests, and correcting it is what gives the
language its central syntactic rule.

When a producer rejects twenty versions and keeps the twenty-first, they have not
merely selected. They have drawn twenty boundaries, each one an act of
individuation: \emph{not that}. The kept version is specified by those rejections
far more precisely than by any positive description, because the positive
descriptions available---``warmer,'' ``nastier,'' ``more forward''---name
directions rather than regions. The rejections name the region.

So the request is not underspecified. It is specified negatively, and the
specification was thrown away, because sessions record survivors and not
exclusions.

Mishima therefore makes the exclusion obligatory. A \texttt{seek} without a
\texttt{not} clause is refused by the parser, not by the type checker
(\Cref{prop:mandatory-not}), on the grounds that a region individuated against
nothing has not been specified at all.

\subsection{Counting commitments, not instructions}

Conventional languages count instructions: a program is expensive because it
executes many. Mishima counts \emph{commitments}, and the distinction reflects
what production actually costs.

Auditioning is cheap. A producer can sweep a filter for an hour, deposit nothing,
and be exactly where they started. Committing is expensive: bouncing a version,
building the next layer on it, resampling it so that everything downstream now
depends on it. The expense is not the time spent but the fact that the decision
is thereafter load-bearing.

Under this model the optimisation target inverts. A program is not good because
it runs fast; it is cheap because it commits rarely. A search that explores an
enormous space of past work while committing three times is inexpensive, however
long its trajectory. \S\ref{sec:cost} makes this precise.

\subsection{Contributions}

\begin{enumerate}[label=(\roman*),leftmargin=2.2em]
\item The obligatory negation and its justification (\S\ref{sec:individuation}).
\item The divergence of refinement cost, and the consequent necessity of
  composing ordinary distinctions (\S\ref{sec:nesting}).
\item The composition law, its saturation dichotomy, and the exact depth required
  to reach a target (\S\ref{sec:nesting}).
\item The triangle result: three mutually supporting rungs are necessary and
  sufficient for robustness to the loss of any one (\S\ref{sec:coherence}).
\item Closure, its strict superiority to a confidence threshold, and the
  two-outcome termination theorem (\S\ref{sec:closure}).
\item Allocation of effort by water-filling, with the knapsack variant for
  probes that are all-or-nothing (\S\ref{sec:allocation}).
\item The complete grammar of \texttt{.mma}, a type system, and static
  diagnostics (\S\ref{sec:grammar}--\S\ref{sec:types}).
\item A prototype and sixteen validating experiments (\S\ref{sec:validation}).
\end{enumerate}

%──────────────────────────────────────────────────────────────────────
\section{Substrate}
\label{sec:substrate}
%──────────────────────────────────────────────────────────────────────

We restate, self-containedly, the minimum structure mishima runs on.

\begin{definition}[Contact graph, medium, cut, residue, floor]
\label{def:substrate}
A \emph{contact graph} is $\Gph=(V,E,\wt)$ with $V$ finite and non-empty, $E$ a
set of unordered pairs, and $\wt\colon E\to\Rpos$; graph terminology follows
\citet{diestel2010graph}. A distinguished vertex $\med$, the \emph{medium}, is
adjacent to every other vertex; the rest are \emph{items}. For $S$ containing an
item and omitting $\med$, the \emph{cut} is
$\cut(S)=\{\{u,v\}\in E: u\in S, v\notin S\}$ and the \emph{residue} is
$\res(S)=\sum_{e\in\cut(S)}\wt(e)$. The \emph{floor} is $\flo=\min_{e\in E}\wt(e)$
and the \emph{total} is $\Om=\sum_{e\in E}\wt(e)$.
\end{definition}

In a production record the items are authored things---a bass patch, a bounce, a
processing move---and an edge weight is the cost of telling two of them apart. The
medium stands for everything not yet distinguished.

\begin{lemma}[Positive floor, positive residue]
\label{lem:floor}
$\flo>0$, and every cut as in \Cref{def:substrate} satisfies $\res(S)\ge\flo>0$.
\end{lemma}

\begin{proof}
$\flo$ is the minimum of a finite non-empty set of positive reals. Since $S$
contains an item $v$ and omits $\med$, and $\{v,\med\}\in E$, the cut is
non-empty and its weight is a sum of terms each at least $\flo$.
\end{proof}

\begin{definition}[Alignment]
\label{def:alignment}
Fix a target $x^{\star}$. The \emph{alignment} of an item $x$ to $x^{\star}$ is
its normalised separation cost,
\[
  a(x,x^{\star}) \;=\; \frac{\aln(x,x^{\star})}{\Om},
  \qquad
  \aln(x,x^{\star}) = \min_{S\ni x,\ x^{\star}\notin S}\res(S),
\]
computable by maximum flow \citep{ford1956maximal,edmonds1972theoretical}.
\end{definition}

Alignment measures distance-to-target as separation cost: an item far from the
target is expensive to separate from it; an item at the target is separated only
at the floor.

\begin{lemma}[Alignment is floor-bounded]
\label{lem:alignfloor}
For all $x$, $a(x,x^{\star})\ge\flo/\Om>0$.
\end{lemma}

\begin{proof}
Immediate from \Cref{lem:floor} applied to the minimising cut.
\end{proof}

\begin{remark}[Uncertainty without a prior]
\label{rem:no-prior}
The quantity ordinarily carried by a probability or a posterior---a number in
$[0,1]$ recording how sure one is---is carried here by $a(\cdot,\cdot)$ and its
distance to $\flo/\Om$. It is derived rather than posited: an exact
minimum-cut weight, bounded below by a positive constant, with no prior to set
and no membership function to tune. Where a standard calculus \emph{posits} an
uncertainty, the floor \emph{forces} one.
\end{remark}

%──────────────────────────────────────────────────────────────────────
\section{Individuation by Negation}
\label{sec:individuation}
%──────────────────────────────────────────────────────────────────────

\begin{definition}[Region, exclusion]
\label{def:region}
A \emph{region} is a set of items together with the cut individuating it from the
remainder. An \emph{exclusion} is a region declared not to be the target.
\end{definition}

\begin{proposition}[A region requires an exclusion]
\label{prop:mandatory-not}
A search specifying a target but no exclusion has not specified a region. The
parser therefore rejects a \texttt{seek} lacking a \texttt{not} clause, and
rejects an empty exclusion list.
\end{proposition}

\begin{proof}
By \Cref{def:region} a region is individuated by a cut, and by
\Cref{def:substrate} a cut is a partition into a side and its complement. A
declaration naming only a side, with the complement left implicit and empty,
determines no partition: every item lies on the named side, the cut is empty, and
by \Cref{lem:floor} no such cut exists over a graph containing the medium. The
declaration therefore fails to denote.
\end{proof}

\begin{remark}[Why this is a parse error and not a type error]
\label{rem:parse-not-type}
The distinction is deliberate. A type error is a claim about the meaning of a
well-formed program; a missing \texttt{not} clause makes the program ill-formed,
because the construct it would inhabit is not a search over a smaller region but
a search over no region at all. Placing the rule in the grammar means the
producer meets it while writing rather than after running, which is where a rule
about specification belongs.
\end{remark}

\begin{example}[A recall, written out]
\label{ex:recall}
The exclusions here are the twenty rejected versions, in the only form that
survived: the properties the producer decided against.

\begin{lstlisting}
seek reese_growl
  not    { thin, undistorted, mono }
  toward { region(that_2019_growl) }
\end{lstlisting}
\end{example}

%──────────────────────────────────────────────────────────────────────
\section{Why Ladders}
\label{sec:nesting}
%──────────────────────────────────────────────────────────────────────

\subsection{A single distinction cannot be sharpened without bound}

\begin{definition}[Refinement cost]
\label{def:refinecost}
Let $\Gamma(t)$ be the cost of realising a separator of thickness $t>0$. We
require $\Gamma$ non-increasing in $t$ with $\Gamma(t)\to\infty$ as
$t\to0^{+}$: making a distinction arbitrarily sharp is arbitrarily expensive.
\end{definition}

The divergence is required of $\Gamma$ rather than derived, and it is not an
arbitrary requirement. Maintaining $N$ mutually distinguishable states inside a
fixed range forces the mean gap between them to $D/N$, so buying a factor of two
in sharpness costs a factor of two in states, and the cost of the mechanism
holding them apart grows at least as fast.

\begin{theorem}[Refinement halts above zero]
\label{thm:refine}
Let the return from thinning a separator be bounded above by the extent $\Lambda$
of the space being resolved. Then the marginal return per unit cost tends to zero
as $t\to0^{+}$, and any agent with a positive threshold on marginal return halts
at a strictly positive thickness.
\end{theorem}

\begin{proof}
Total return is bounded by $\Lambda$ regardless of $t$, so the return available
from any further thinning is at most $\Lambda$ minus what has been obtained,
which is non-increasing in the thinning already performed. Cost, by
\Cref{def:refinecost}, diverges. The ratio of a bounded non-increasing numerator
to a divergent denominator tends to zero, so it falls below any positive
threshold at some $t^{\ast}>0$.
\end{proof}

The experiment \texttt{mishima\_refinement\_diverges} exhibits the divergence
directly: thicknesses $10^{-1}$ through $10^{-5}$ cost $10$ through $10^{5}$, a
ratio of $10^{4}$ across the sweep.

\subsection{Rungs and ladders}

\begin{definition}[Rung, power, ladder]
\label{def:rung}
A \emph{rung} is a move that closes some of the distance from the current item to
the target. Its \emph{power} is
\[
  \pot(\gamma,x,x^{\star})
  \;=\;
  \frac{a(x,x^{\star}) - a(\gamma(x),x^{\star})}
       {a(x,x^{\star}) - \flo/\Om}\;\in\;[0,1],
\]
inert at $\pot=0$ and absolute at $\pot=1$. A \emph{ladder} is a finite sequence
of rungs composed in series, written $\gamma_1 \gg \cdots \gg \gamma_n$.
\end{definition}

\begin{theorem}[Composition is multiplicative in the residual]
\label{thm:multiplicative}
A ladder of rungs with powers $\pot_1,\dots,\pot_n$ has composite power
\[
  \pot_{\mathrm{total}} \;=\; 1-\prod_{i=1}^{n}\left(1-\pot_i\right).
\]
\end{theorem}

\begin{proof}
By \Cref{def:rung} a rung of power $\pot_i$ leaves a fraction $1-\pot_i$ of the
above-floor gap unclosed. Applying rungs in series multiplies the surviving
fractions, so the fraction remaining after $n$ rungs is $\prod_i(1-\pot_i)$ and
the fraction closed is its complement.
\end{proof}

The experiment \texttt{mishima\_composition\_law} checks the implementation
against a direct product over four cases, with greatest absolute difference below
$10^{-12}$.

\begin{corollary}[Saturation dichotomy]
\label{cor:saturation}
The residual gap $\prod_i(1-\pot_i)$ tends to zero if and only if
$\sum_i \pot_i$ diverges.
\end{corollary}

\begin{proof}
Taking logarithms, $\log\prod_i(1-\pot_i)=\sum_i\log(1-\pot_i)$. For
$\pot_i\in[0,1)$ we have $-\log(1-\pot_i)\ge \pot_i$ and, for $\pot_i$ bounded
away from one, $-\log(1-\pot_i)\le C\pot_i$ for a constant $C$. The sum therefore
diverges to $-\infty$---equivalently the product tends to zero---exactly when
$\sum_i\pot_i$ diverges. This is the Borel--Cantelli dichotomy in the form given
by \citet{borel1909probabilites}.
\end{proof}

\begin{table}[H]
\centering
\small
\caption{Saturation over $4000$ rungs
(\texttt{mishima\_saturation\_dichotomy}).}
\label{tab:saturation}
\begin{tabular}{@{}lrr@{}}
\toprule
Power sequence & $\sum\pot_i$ & Residual gap \\
\midrule
Harmonic, $\pot_i = 1/(i{+}2)$   & divergent  & $2.50\times10^{-4}$ \\
Geometric, $\pot_i = 2^{-(i+1)}$ & convergent & $2.89\times10^{-1}$ \\
\bottomrule
\end{tabular}
\end{table}

\begin{corollary}[Required depth]
\label{cor:depth}
To reach composite power $\pot^{\star}<1$ using rungs of power $\pot$ requires
\[
  n^{\star} \;=\;
  \left\lceil \frac{\log(1-\pot^{\star})}{\log(1-\pot)} \right\rceil
\]
rungs, and this is exact rather than indicative.
\end{corollary}

\begin{proof}
By \Cref{thm:multiplicative} with equal powers the composite is $1-(1-\pot)^n$;
solving $1-(1-\pot)^n\ge\pot^{\star}$ for $n$ gives the stated ceiling.
\end{proof}

The experiment \texttt{mishima\_depth\_closed\_form} compares the closed form
against incrementing $n$ until the composite clears the target, over twenty
combinations of target and per-rung power, and finds agreement in every case.

\begin{remark}[The ladder is the computation]
\label{rem:ladder-is}
\Cref{thm:refine} and \Cref{thm:multiplicative} together say that fine resolution
is reached by paying many ordinary costs in series rather than one divergent cost
at a single boundary. This is not scaffolding for the search; it is the search.
The production analogue is exact and familiar: a sound that survives ten passes
of resampling and mangling is not the product of one perfect operation but of ten
ordinary ones, and no amount of care applied to any single stage substitutes for
the sequence.
\end{remark}

\begin{corollary}[Diversify, against the mean]
\label{cor:diversify}
A ladder of distinct rungs attains strictly greater composite power than
repetition of any one of its rungs whose power lies below the ladder's mean.
\end{corollary}

\begin{proof}
By \Cref{thm:multiplicative} the composite is determined by
$\prod_i(1-\pot_i)$, and by the arithmetic--geometric mean inequality applied to
the factors $1-\pot_i$, replacing the set by $n$ copies of a member with
$\pot<\bar\pot$ increases the product and hence decreases the composite.
\end{proof}

\begin{remark}[The stronger claim is false]
\label{rem:diversify-honest}
It is tempting to state \Cref{cor:diversify} as ``distinct rungs beat any
repetition,'' and that is not true. Repeating the \emph{strongest} rung of a set
can beat the mixed ladder: for powers $\{0.45,0.30,0.55\}$ the ladder attains
$0.8267$ while three copies of $0.55$ attain $0.9089$. The experiment
\texttt{mishima\_diversify} records this counterexample rather than suppressing
it. What the corollary licenses is the comparison against the mean, which is
weaker and true. The practical reading is that a ladder buys robustness rather
than raw power---for which see \Cref{thm:triangle}.
\end{remark}

%──────────────────────────────────────────────────────────────────────
\section{Coherence}
\label{sec:coherence}
%──────────────────────────────────────────────────────────────────────

\begin{definition}[Support relation]
\label{def:support}
Rung $\gamma_i$ \emph{supports} $\gamma_j$ if the region $\gamma_i$ reaches is
one from which $\gamma_j$'s power is realisable. The support relation is a
directed graph on the rungs of a ladder.
\end{definition}

\begin{theorem}[Triangle]
\label{thm:triangle}
A ladder survives the loss of any single rung if and only if its support relation
contains a directed cycle of length at least three. In particular a chain is
fragile and a mutually supporting pair is not sufficient.
\end{theorem}

\begin{proof}
($\Leftarrow$) If a directed cycle of length $k\ge3$ is present, removing one
rung leaves a directed path through the remaining $k-1\ge2$, so every surviving
rung is still reached and the ladder is still traversable.
($\Rightarrow$) Suppose no cycle of length $\ge3$ exists. Then the support
relation is either acyclic or contains only $2$-cycles. If acyclic, it has a
topological order and removing the first rung disconnects everything downstream
of it. If it contains a $2$-cycle $\{\gamma_i,\gamma_j\}$ and no longer cycle,
removing $\gamma_i$ leaves $\gamma_j$ supported only by the rung just removed.
Either way some single removal breaks the ladder.
\end{proof}

\begin{remark}[Two agreeing is not a check]
\label{rem:two-agree}
A $2$-cycle is mutual definition between exactly two rungs with no independent
third party. In a recall this is the situation where a spectral measurement and a
stored annotation agree because the annotation was written from the measurement;
their agreement is not evidence. The requirement of three is what makes agreement
informative, and it is checkable statically from the signs of the support
relation alone---no magnitudes are needed.
\end{remark}

The experiment \texttt{mishima\_triangle} confirms that a chain and a pair are
both fragile and a three-cycle is not.

%──────────────────────────────────────────────────────────────────────
\section{Closure and Decline}
\label{sec:closure}
%──────────────────────────────────────────────────────────────────────

A search over a production record cannot terminate by matching an expectation,
because if an expectation were available the search would be unnecessary. We give
a criterion that requires none.

\begin{definition}[Reached set]
\label{def:reached}
At a given stage the \emph{reached set} $\Reached$ of a search is the collection
of regions attainable by some admissible ladder from the seed using the rungs
invoked so far, partitioned into classes under indistinguishability at the
endpoints.
\end{definition}

\begin{definition}[Closure]
\label{def:closure}
A search is \emph{closed} at a stage if, for every rung available but not yet
invoked, extending by that rung adds no new class to $\Reached$.
\end{definition}

\begin{theorem}[Closure is strictly stronger than a threshold]
\label{thm:closure-stronger}
There exist records and thresholds $\theta\in(0,1)$ such that a search's terminal
alignment satisfies $\theta$ while the search is not closed: an uninvoked rung is
available whose ladder reaches a region in a distinct class.
\end{theorem}

\begin{proof}
Construct a record with two disjoint clusters $A$ and $B$, each internally dense
and each joined to the medium by a single rung, $\gamma_A$ and $\gamma_B$, of
equal weight. A ladder using $\gamma_A$ alone terminates at a representative
$x^{\star}_A\in A$, whose alignment to itself is $\flo/\Om$ by
\Cref{lem:alignfloor} and therefore satisfies any threshold. Yet $\gamma_B$,
uninvoked, reaches $x^{\star}_B\in B$; as vertex sets with independent internal
structure their minimum cuts against the medium differ, so the two lie in
distinct classes. The threshold is met while closure fails.
\end{proof}

\begin{remark}[The failure this names]
\label{rem:why-threshold}
\Cref{thm:closure-stronger} formalises a failure familiar to anyone who has
searched their own archive. One line of enquiry reports high confidence purely
because it is internally consistent---as any single cluster is, having only
internal contacts and one to the medium---while an independent second line would
reach an entirely different answer. Closure requires that every available rung be
checked to land in an already-reached class before stopping. It is not that the
current answer looks confident; it is that the space of answers has stopped
producing new destinations.
\end{remark}

\begin{table}[H]
\centering
\small
\caption{Closure against a fixed confidence threshold $\theta=0.9$, over five
registries (\texttt{mishima\_closure\_beats\_threshold}).}
\label{tab:closure}
\begin{tabular}{@{}rrrl@{}}
\toprule
Registry & Threshold met & Closure at & Outcome \\
\midrule
$3$ & after $1$ probe & $3$ probes & declined \\
$4$ & after $1$ probe & $4$ probes & declined \\
$5$ & after $1$ probe & $5$ probes & declined \\
$6$ & after $1$ probe & $6$ probes & declined \\
$8$ & after $1$ probe & $8$ probes & declined \\
\bottomrule
\end{tabular}
\end{table}

\begin{theorem}[Convergent closure or honest decline]
\label{thm:decline}
Every search over a finite rung registry terminates in exactly one of two states:
\begin{enumerate}[label=(\roman*),leftmargin=2.2em]
\item \textbf{Convergent closure}: $\Reached$ stabilises to a single class, whose
  representative is reported;
\item \textbf{Contested closure (honest decline)}: $\Reached$ stabilises with
  more than one class, and the search reports that no single sufficient answer
  exists, together with the classes found.
\end{enumerate}
No search over a finite registry runs forever without reaching one of these.
\end{theorem}

\begin{proof}
The registry is finite, so the set of ladders from the seed is finite and
$\Reached$---a partition of a finite set of reachable regions---is finite. It
grows only by adding classes and can do so at most finitely often; once each
remaining rung has been tried against every present class, no rung adds a class
and \Cref{def:closure} is satisfied by exhaustion in finitely many steps. If the
stabilised $\Reached$ has one class, (i) obtains; if more, no class is licensed
over the others, since each is internally indistinguishable while they disagree
mutually, and (ii) obtains. The cases are exhaustive and mutually exclusive.
\end{proof}

\begin{remark}[Decline is not failure]
\label{rem:decline-value}
A contested closure reports a fact about the record: several irreconcilable
regions were reached and nothing in the available rungs distinguishes them.
Reporting one and suppressing the others would be a stronger claim than the
search supports. For a producer this is the more useful outcome: told that a
sound could be either of two things and named the measurement that would separate
them, they know what to listen for. A threshold-stopping procedure discards
exactly this information.
\end{remark}

The experiment \texttt{mishima\_decline\_is\_typed} exhibits both outcomes: a
registry reaching two classes terminates contested carrying both, while one
reaching a single class converges.

%──────────────────────────────────────────────────────────────────────
\section{Allocating Effort}
\label{sec:allocation}
%──────────────────────────────────────────────────────────────────────

A search has finite effort and many rungs it might invoke.

\begin{definition}[Yield]
\label{def:yield}
Rung $i$ given effort $a_i\ge0$ returns yield $y_i(a_i)$, non-decreasing and
concave, with $\sum_i a_i = E$ the total effort available.
\end{definition}

\begin{theorem}[Allocation is water-filling]
\label{thm:waterfill}
The allocation maximising total yield satisfies, for a single scalar price
$\lambda\ge0$,
\[
  y_i'(a_i^{\star}) = \lambda \ \text{ where } a_i^{\star}>0,
  \qquad
  y_i'(0) \le \lambda \ \text{ where } a_i^{\star}=0,
\]
with $\sum_i a_i^{\star}=E$.
\end{theorem}

\begin{proof}
Maximising a sum of concave functions under a linear budget constraint is a
convex program; the stated conditions are its Karush--Kuhn--Tucker conditions
\citep{boyd2004convex}, with $\lambda$ the multiplier on the budget. Concavity
makes them sufficient as well as necessary.
\end{proof}

\begin{corollary}[Effort is refused, not rationed]
\label{cor:engaged}
Effort goes to exactly those rungs whose marginal yield at zero exceeds the
price, and none to the rest. The price is a single readable scalar recording how
loaded the search is.
\end{corollary}

The experiment \texttt{mishima\_water\_filling} allocates a budget of $3$ across
three rungs, spends the budget to within $10^{-3}$, equalises the marginal yields
of the served rungs to within $10^{-3}$, and gives nothing to the weakest.

\begin{remark}[When concavity fails]
\label{rem:knapsack}
Concavity is the one premise here that is a claim about the rungs rather than
about the structure allocating among them, and audio rungs frequently violate it.
A stem separator returns nothing usable below some effort and a great deal above
it: a threshold, not a concave curve. Where that holds, the correct formulation
is a knapsack, admitting rungs in decreasing order of
\[
  \varrho_i \;=\; \frac{1}{c_i}\,\log\frac{\Om}{\Om-\flo_i},
\]
which the experiment \texttt{mishima\_knapsack\_priority} confirms ranks a cheap
strong rung above a dear strong one and a cheap weak one.
\end{remark}

%──────────────────────────────────────────────────────────────────────
\section{The Cost Model}
\label{sec:cost}
%──────────────────────────────────────────────────────────────────────

\begin{definition}[Probe and commit]
\label{def:probecommit}
To \emph{probe} is to traverse the record and arrive at a region: a fact about
the trajectory, drawing no boundary. To \emph{commit} is to individuate a region
so that it can be named, stored and built upon: an act depositing residue at
least $\flo$.
\end{definition}

\begin{theorem}[Cost is the commitment count]
\label{thm:cost}
The cost of a mishima computation is the number of commitments it makes, and is
independent of the number of probes it performs.
\end{theorem}

\begin{proof}
By \Cref{def:probecommit} a probe draws no cut and so deposits no residue; by
\Cref{lem:floor} each commitment deposits at least $\flo$. Total residue is
therefore $\sum$ over commitments alone, in which no term indexed by a probe
appears.
\end{proof}

\begin{remark}[What this licenses, and what it does not]
\label{rem:cost-honest}
\Cref{thm:cost} is a statement about residue, not about wall-clock time. A search
performing a million probes takes longer than one performing ten, and an
implementation must still care. What the theorem says is that the long search has
not made the record any less sharp, whereas ten careless commitments have. The
engineering question it recommends is therefore not how to search faster but
where to commit, and the answer to the second question does not follow from the
first.
\end{remark}

\begin{corollary}[Auditioning is free]
\label{cor:audition}
A producer who sweeps a parameter for an hour and commits nothing has left the
record exactly as they found it.
\end{corollary}

%──────────────────────────────────────────────────────────────────────
\section{Grammar}
\label{sec:grammar}
%──────────────────────────────────────────────────────────────────────

\subsection{Lexical structure}

Source is UTF-8 with extension \texttt{.mma}. Comments run from \texttt{--} to
end of line. The language is not layout-sensitive.

\begin{table}[H]
\centering
\small
\caption{Token classes.}
\label{tab:tokens}
\begin{tabular}{@{}ll@{}}
\toprule
Class & Definition \\
\midrule
\textsc{Ident} & \texttt{[A-Za-z\_][A-Za-z0-9\_]*}, not a keyword \\
\textsc{Num}   & \texttt{-?[0-9]+(.[0-9]+)?([eE][+-]?[0-9]+)?} \\
\textsc{Str}   & double-quoted, standard escapes \\
\textsc{Op}    & \texttt{>> := <= >= == != \{ \} ( ) , : . < > \#} \\
\bottomrule
\end{tabular}
\end{table}

Keywords: \texttt{floor}, \texttt{module}, \texttt{probe}, \texttt{commit},
\texttt{seek}, \texttt{not}, \texttt{toward}, \texttt{via}, \texttt{until},
\texttt{closure}, \texttt{converge}, \texttt{otherwise}, \texttt{decline},
\texttt{yield}, \texttt{ladder}, \texttt{rung}, \texttt{at}, \texttt{observe},
\texttt{assert}, \texttt{emit}, \texttt{let}, \texttt{as}, \texttt{when}.

A numeric literal may carry an explicit resolution by the suffix
\texttt{value\#floor}, for example \texttt{6.0\#0.5}. A literal without one
inherits the ambient floor. \emph{There is no way to write a literal of zero
resolution}: the lexer rejects a non-positive floor suffix, which the experiment
\texttt{mishima\_lexer\_floor\_suffix} confirms over three attempts.

\subsection{Context-free grammar}

\begin{lstlisting}[caption={Mishima grammar (EBNF).},label={lst:ebnf}]
program   ::= { decl }
decl      ::= floorDecl | moduleDecl | stmt
floorDecl ::= "floor" num
moduleDecl::= "module" ident "{" { decl } "}"

stmt      ::= bind | seek | observe | assert | expr
bind      ::= [ "let" ] ident ":=" expr

seek      ::= "seek" ident
                "not"    "{" exclList "}"
                "toward" "{" expr "}"
              [ "via"    "{" ladder "}" ]
                "until"  admit
              [ "otherwise" "decline" ]
                "yield"  ident

ladder    ::= rungExpr { ">>" rungExpr }
rungExpr  ::= "rung" ident [ "at" num ]
            | ident "(" [ argList ] ")"

admit     ::= "closure" | "converge" | cond

expr      ::= "probe"  "(" expr ")"
            | "commit" "(" expr ")"
            | "ladder" "{" ladder "}"
            | ident | num | str | "(" expr ")"

observe   ::= "observe" expr [ "as" ident ]
assert    ::= "assert" cond [ "emit" str ]
exclList  ::= expr { "," expr }
argList   ::= arg { "," arg }
arg       ::= [ ident ":" ] expr
cond      ::= expr relop expr
relop     ::= "<" | ">" | "<=" | ">=" | "==" | "!="
num       ::= NUM [ "#" NUM ]
\end{lstlisting}

The \texttt{not} clause is mandatory: a \texttt{seek} lacking one is rejected at
parse time (\Cref{prop:mandatory-not}), as is an empty exclusion list. The
\texttt{via} clause is optional; when omitted the checker derives a ladder from
the registry until the composite power meets the target
(\Cref{thm:multiplicative}).

\subsection{Worked program}

\begin{lstlisting}[caption={A complete mishima program.},label={lst:example}]
-- recall.mma
floor 0.02                 -- nothing here resolves finer

seek reese_growl
  not    { thin, undistorted, mono }
  toward { region(that_2019_growl) }
  via    { rung spectral   at 0.45
        >> rung annotation at 0.30
        >> rung model      at 0.55 }
  until  closure
  otherwise decline
  yield  found

observe found as result
assert result.residue >= 0.02
  emit "residue fell below the declared floor"
\end{lstlisting}

The ladder has composite power
$1-(1-0.45)(1-0.30)(1-0.55)=1-0.55\times0.70\times0.45=0.82675$, closing about
$82.7\%$ of the above-floor gap. Three rungs of differing kind are used rather
than three of one, so that the support structure can form a triangle
(\Cref{thm:triangle}); the experiment
\texttt{mishima\_parses\_worked\_example} recovers exactly this figure from the
source.

\begin{remark}[The three rungs are deliberately unlike]
\label{rem:three-unlike}
\texttt{spectral} measures the material, \texttt{annotation} recalls what was
written at the time, and \texttt{model} reads the accumulated graph. They fail in
unrelated ways, which is what \Cref{rem:two-agree} requires. Three prompts to one
model would be one rung, not three.
\end{remark}

%──────────────────────────────────────────────────────────────────────
\section{Type System}
\label{sec:types}
%──────────────────────────────────────────────────────────────────────

\begin{definition}[Types]
\label{def:types}
\[
  \tau ::= \textsf{Region} \mid \textsf{Rung} \mid \textsf{Ladder}
        \mid \textsf{Path} \mid \textsf{Scalar} \mid \textsf{Result}
\]
where $\textsf{Result}$ is the sum
$\textsf{Converged}(\textsf{Path})+\textsf{Declined}(\{\textsf{Region}\})$.
\end{definition}

\begin{definition}[Judgement]
\label{def:judgement}
The typing judgement, in the style of \citet{pierce2002types}, is
\[
  \Gamma \vdash e : \tau \,@\, \flo \,/\, \pot,
\]
read: in context $\Gamma$, expression $e$ has type $\tau$, is individuated at
resolution $\flo>0$, and carries power $\pot\in[0,1]$. Scalars carry $\pot=0$.
\end{definition}

\subsection{Rules}

\[
\frac{\Gamma \vdash n \quad \flo > 0}
     {\Gamma \vdash n\#\flo : \textsf{Scalar} \,@\, \flo \,/\, 0}
\ \textsc{(Lit)}
\]

\[
\frac{\Gamma \vdash e : \textsf{Region} \,@\, \flo \,/\, \pot}
     {\Gamma \vdash \texttt{probe}(e) : \textsf{Region} \,@\, \flo \,/\, 0}
\ \textsc{(Probe)}
\]

\[
\frac{\Gamma \vdash e : \textsf{Region} \,@\, \flo \,/\, \pot \quad \flo > 0}
     {\Gamma \vdash \texttt{commit}(e) : \textsf{Path} \,@\, \flo \,/\, \pot}
\ \textsc{(Commit)}
\]

\[
\frac{\Gamma \vdash \gamma_i : \textsf{Rung} \,@\, \flo \,/\, \pot_i}
     {\Gamma \vdash \gamma_1 \gg \cdots \gg \gamma_n :
      \textsf{Ladder} \,@\, \flo \,/\, 1-\prod_i(1-\pot_i)}
\ \textsc{(Comp)}
\]

\[
\frac{\begin{array}{c}
       \Gamma \vdash N : \{\textsf{Region}\} \quad N \neq \emptyset \\
       \Gamma \vdash t : \textsf{Region} \quad
       \Gamma \vdash L : \textsf{Ladder} \,@\, \flo \,/\, \pot
      \end{array}}
     {\Gamma \vdash \texttt{seek}\ x\ \texttt{not}\ N\ \texttt{toward}\ t\
       \texttt{via}\ L : \textsf{Result} \,@\, \flo \,/\, \pot}
\ \textsc{(Seek)}
\]

Rule \textsc{(Probe)} is the type-level statement that probing closes no distance
and deposits nothing. Rule \textsc{(Commit)} carries the power through and yields
a $\textsf{Path}$, the type of things that have cost something. The side
condition $\flo>0$ appears in every value-introducing rule, and the premise
$N\neq\emptyset$ enforces \Cref{prop:mandatory-not} at the type level as well as
in the grammar.

\begin{theorem}[No zero-resolution value]
\label{thm:no-zero}
There is no well-typed mishima expression of a region-denoting type with
resolution $0$.
\end{theorem}

\begin{proof}
Every value-introducing rule carries the side condition $\flo>0$, and no rule
lowers the resolution annotation of an already-typed value: \textsc{(Comp)}
alters only the power, and \textsc{(Probe)} and \textsc{(Commit)} carry $\flo$
through unchanged. By induction on the derivation every region-denoting value
satisfies $\flo>0$.
\end{proof}

\subsection{Static diagnostics}

\begin{proposition}[Coherence diagnostic]
\label{prop:cohdiag}
A \texttt{via} ladder of length one or two declaring \texttt{until closure} is
rejected at compile time.
\end{proposition}

\begin{proof}
By \Cref{thm:triangle} a robust support structure requires a cycle of length at
least three, which a ladder of fewer than three rungs cannot contain. The check
needs only the rung count and the signs of the support relation, both available
statically.
\end{proof}

\begin{proposition}[Saturation diagnostic]
\label{prop:satdiag}
If a \texttt{seek} declares a target requiring composite power $\pot^{\star}$ and
its ladder attains less, the program is rejected with the attainable power and
the required depth named.
\end{proposition}

\begin{proof}
Composite power is computable from the rung powers, which \textsc{(Comp)} records
in the type, and the depth follows from \Cref{cor:depth}. Both comparisons are
therefore decidable statically.
\end{proof}

\begin{principle}[Every refusal carries a remedy]
\label{prin:remedy}
A diagnostic states the rule that fired, what is wrong, and what to do about it.
A refusal with no way forward is a dead end, and the checker emits none.
\end{principle}

The experiment \texttt{mishima\_coherence\_diagnostic} confirms that ladders of
length one and two are refused while a ladder of three is accepted, and
\texttt{mishima\_floor\_negotiation} confirms that every diagnostic carries a
non-empty remedy.

%──────────────────────────────────────────────────────────────────────
\section{Operational Semantics}
\label{sec:semantics}
%──────────────────────────────────────────────────────────────────────

A configuration is $\langle e, \Sigma, \rec\rangle$: an expression, an evidence
store recording the regions committed so far, and a monotone counter $\rec$
recording commitments.

\[
\frac{}{\langle \texttt{probe}(e), \Sigma, \rec\rangle \to
        \langle v_{\textsf{Region}}, \Sigma, \rec\rangle}
\ \textsc{(E-Probe)}
\]

\[
\frac{}{\langle \texttt{commit}(e), \Sigma, \rec\rangle \to
        \langle v_{\textsf{Path}}, \Sigma\oplus e, \rec+1\rangle}
\ \textsc{(E-Commit)}
\]

\[
\frac{\Reached \text{ stable},\ |\Reached| = 1}
     {\langle \texttt{seek}\dots, \Sigma, \rec\rangle \to
      \langle \textsf{Converged}(p), \Sigma, \rec + k\rangle}
\ \textsc{(E-Close)}
\]

\[
\frac{\Reached \text{ stable},\ |\Reached| > 1}
     {\langle \texttt{seek}\dots, \Sigma, \rec\rangle \to
      \langle \textsf{Declined}(\Reached), \Sigma, \rec + k\rangle}
\ \textsc{(E-Decline)}
\]

Note that \textsc{(E-Probe)} leaves both $\Sigma$ and $\rec$ unchanged while
\textsc{(E-Commit)} advances both: this is \Cref{thm:cost} as a reduction rule.

\begin{theorem}[Progress and preservation]
\label{thm:soundness}
If $\Gamma\vdash e:\tau\,@\,\flo\,/\,\pot$ then either $e$ is a value or there is
$\langle e',\Sigma',\rec'\rangle$ with
$\langle e,\Sigma,\rec\rangle\to\langle e',\Sigma',\rec'\rangle$; and in the
latter case $\Gamma\vdash e':\tau\,@\,\flo'\,/\,\pot'$ with $\flo'\ge\flo$ and
$\rec'\ge\rec$.
\end{theorem}

\begin{proof}
By induction on the typing derivation, in the syntactic style of
\citet{wright1994syntactic}. Progress: each non-value form is the subject of
exactly one reduction rule whose premises are discharged by the typing premises,
and by \Cref{thm:decline} the \texttt{seek} form always reduces, since a finite
registry stabilises. Preservation: each rule's conclusion has the type its
premise assigns, resolution is carried through unchanged or increased, and
$\rec$ is incremented only by \textsc{(E-Commit)}.
\end{proof}

\begin{corollary}[No stuck state]
\label{cor:nostuck}
A well-typed mishima program does not get stuck, and every value it produces
carries a positive resolution.
\end{corollary}

%──────────────────────────────────────────────────────────────────────
\section{On Models as Rungs}
\label{sec:models}
%──────────────────────────────────────────────────────────────────────

A rung may be a spectral measurement, a stored annotation, a deterministic rule,
or an inference by a trained model. Nothing in \Cref{def:rung} or in any result
above distinguishes these cases, because no definition refers to anything but a
rung's action on items and the alignment it closes.

\begin{proposition}[Rung kind is invisible]
\label{prop:kind-invisible}
Two rungs of equal power and cost are interchangeable in every result of this
paper.
\end{proposition}

\begin{proof}
Immediate: power and cost are the only rung attributes appearing in
\Cref{thm:multiplicative}, \Cref{thm:waterfill}, \Cref{thm:triangle} or
\Cref{thm:decline}.
\end{proof}

The experiment \texttt{mishima\_probe\_kind\_invisible} substitutes a model
inference for a spectral measurement of equal power and obtains an identical
outcome.

\begin{remark}[What a model may and may not do here]
\label{rem:model-bounded}
\Cref{prop:kind-invisible} is a licence and a limit. A model may serve as a rung,
and the language will neither privilege nor discount it. What the language does
not offer is any way for a model to adjudicate: there is no construct in which a
model's output overrides a contested closure, because by \Cref{thm:decline} the
contested outcome reports a fact about the record---that several irreconcilable
regions were reached---and no rung can make that fact untrue. A model that
disagrees with the other rungs produces a decline naming the disagreement, which
is the correct report.
\end{remark}

\begin{remark}[Classification of rungs is scheduling, not semantics]
\label{rem:kind-scheduling}
Distinguishing rungs by kind remains useful for allocation, since expected cost
and latency differ and \Cref{thm:waterfill} consumes both. It carries no weight
in the theory.
\end{remark}

%──────────────────────────────────────────────────────────────────────
\section{Prototype and Validation}
\label{sec:validation}
%──────────────────────────────────────────────────────────────────────

A prototype accompanies this paper: a lexer, an LL(1) recursive-descent parser, a
checker implementing the static diagnostics, and the ladder and closure
machinery, in \texttt{hkcore.py} with experiments in \texttt{exp\_mishima.py}. It
uses the Python standard library only. The full suite across all three papers
reports \textbf{40 run, 40 passed, 0 failed} in $0.16$\,s, of which sixteen
experiments belong to this paper. Each fixes its seed; re-running reproduces
every record bit-identically.

\begin{table}[H]
\centering
\small
\caption{Experiments for this paper.}
\label{tab:experiments}
\begin{tabular}{@{}lp{4.2cm}@{}}
\toprule
Experiment & Establishes \\
\midrule
\texttt{lexer\_floor\_suffix} & no zero-resolution literal; 3 refusals \\
\texttt{parses\_worked\_example} & \Cref{lst:example} yields $\pot=0.8267$ \\
\texttt{mandatory\_not} & \Cref{prop:mandatory-not} at parse time \\
\texttt{composition\_law} & \Cref{thm:multiplicative}; error $<10^{-12}$ \\
\texttt{diversify} & \Cref{cor:diversify} and its counterexample \\
\texttt{saturation\_dichotomy} & \Cref{cor:saturation}; 3 orders apart \\
\texttt{depth\_closed\_form} & \Cref{cor:depth}; 20/20 agree \\
\texttt{refinement\_diverges} & \Cref{thm:refine}; cost ratio $10^{4}$ \\
\texttt{triangle} & \Cref{thm:triangle}; chain and pair fragile \\
\texttt{coherence\_diagnostic} & \Cref{prop:cohdiag}; 1 and 2 refused \\
\texttt{closure\_beats\_threshold} & \Cref{thm:closure-stronger}; 5 registries \\
\texttt{decline\_is\_typed} & \Cref{thm:decline}; both outcomes \\
\texttt{water\_filling} & \Cref{thm:waterfill}; spread $<10^{-3}$ \\
\texttt{knapsack\_priority} & \Cref{rem:knapsack}; ranking correct \\
\texttt{probe\_kind\_invisible} & \Cref{prop:kind-invisible} \\
\texttt{floor\_negotiation} & \Cref{prin:remedy}; every remedy non-empty \\
\bottomrule
\end{tabular}
\end{table}

\begin{remark}[A negative result, recorded]
\label{rem:negative}
\texttt{mishima\_diversify} began as a check of the claim that distinct rungs beat
every repetition. That claim is false, and the experiment was rewritten to record
the counterexample rather than the claim. We state this because a suite that
never contradicts its paper is not evidence of anything; see
\Cref{rem:diversify-honest}.
\end{remark}

%──────────────────────────────────────────────────────────────────────
\section{Related Notions}
\label{sec:related}
%──────────────────────────────────────────────────────────────────────

\emph{Retrieval.} Similarity search over audio embeddings
\citep{casey2008content} answers ``what resembles this?'' Mishima answers a
different question---``what did I do, and what was essential to it?''---and the
difference is not one of quality. A retrieval system ranks candidates by
proximity; \Cref{thm:decline} declines to rank when the record does not license a
single answer.

\emph{Query languages over provenance.} Provenance query languages
\citep{davidson2008provenance} share the concern with recovering how a result
arose. The departure is the termination criterion: closure quantifies over
uninvoked rungs (\Cref{def:closure}) rather than over recorded facts, and so can
report that the record has stopped producing new destinations.

\emph{Anytime algorithms.} The literature on anytime computation
\citep{zilberstein1996using} supplies the idea that a computation should be
interruptible with a usable partial answer. \Cref{thm:decline} sharpens the
stopping rule: the correct condition is non-resolution against the floor rather
than budget exhaustion, and a search that has not resolved should decline rather
than report its best guess.

\emph{Convex allocation.} \Cref{thm:waterfill} is the standard water-filling
solution \citep{boyd2004convex}, here applied to attention across rungs rather
than to power across channels.

%──────────────────────────────────────────────────────────────────────
\section{Limits}
\label{sec:limits}
%──────────────────────────────────────────────────────────────────────

\emph{The record must exist.} Mishima queries a graph of authored decisions and
cannot construct one retrospectively. A producer with ten years of bounces and no
record has nothing for it to read, and \S\ref{sec:intro} argues the reconstruction
is not available. The language is worth adopting only prospectively.

\emph{Powers must be supplied.} \Cref{def:rung} takes a rung's power as given.
Calibrating what a spectral measurement or an annotation actually closes, in a
particular archive, is a measurement problem this paper does not solve, and a
badly calibrated registry gives a well-typed program that reaches the wrong
place.

\emph{Set semantics, not probabilistic.} A rung that is wrong beyond its declared
resolution can empty the reached set, and the language reports this as a
degenerate decline rather than down-weighting the offending rung. This is a
deliberate choice (\Cref{rem:no-prior}) with a real cost: there is no graceful
degradation under a systematically miscalibrated rung.

\emph{Nothing here evaluates music.} The language finds and relates; it holds no
expectation and issues no verdict. A search that declines has said something
true about the record, not something about the material's quality.

%──────────────────────────────────────────────────────────────────────
\section{Conclusion}
\label{sec:conclusion}
%──────────────────────────────────────────────────────────────────────

We have specified a language whose central syntactic rule is that a search must
say what it excludes, on the grounds that a region individuated against nothing
is not a region (\Cref{prop:mandatory-not}); whose cost model counts commitments
rather than instructions, so that exploration is free and the engineering
question becomes where to commit (\Cref{thm:cost}); whose resolution is reached by
composing ordinary distinctions rather than sharpening one, because sharpening
diverges (\Cref{thm:refine}, \Cref{thm:multiplicative}); and which terminates on
closure rather than confidence, reporting a typed decline when the record
supports no single answer (\Cref{thm:closure-stronger},
\Cref{thm:decline}).

What this offers a producer is the ability to ask, of their own accumulated work,
a question they cannot presently ask of anything: not \emph{what resembles this}
but \emph{what did I do, what was essential to it, and where does the evidence
disagree with itself}. The last of those is the one a threshold would have
thrown away, and we think it is frequently the most useful.

\bibliographystyle{plainnat}
\bibliography{references}

\end{document}
